\chapter{Séries temporelles linéaires}\label{ch:qm:linear-time-series}

Le spread entre les rendements des bons du Trésor américain à dix ans et à deux ans, le 2s10s négocié par un steepener, présente une autocorrélation quotidienne de
0.9987 sur cinquante ans. Une régression du spread de demain sur celui d'aujourd'hui donne une pente de 0.99875 avec une \omterm{def:qm:estimation:estimator}{erreur type} de 0.00045~: une
statistique $t$ de $-2.79$ contre une pente égale à un, une $p$-valeur unilatérale de 0.003 selon la table normale. Le spread semble revenir à la moyenne,
avec une \omterm{def:qm:stochastic-differential-equations:ou}{demi-vie} de 552 jours de bourse. Mais si le spread avait une \omterm{def:qm:linear-time-series:unitroot}{racine unitaire}, la table normale serait la mauvaise~: la bonne table place la
valeur critique à 5~\%, soit $-2.86$, et la preuve est insuffisante. Ce chapitre traite des modèles qui décrivent ces séries et des tests qui
permettent de les distinguer~: la stationnarité et la \omterm{def:qm:linear-time-series:acf}{fonction d'autocorrélation}, les modèles autorégressifs et à moyenne mobile, la vision spectrale, les \omterm{def:qm:linear-time-series:unitroot}{racines unitaires}
et les régressions qu'elles rendent fallacieuses, ainsi que la \omterm{def:qm:linear-time-series:longmem}{mémoire longue}.

\section{Stationnarité et autocorrélation}

\begin{definition}[Processus stationnaire, stationnarité faible, bruit blanc]\label{def:qm:linear-time-series:stationary}
Un processus $(X_t)$ est un \emph{processus stationnaire}\index{processus stationnaire} (au sens strict) si la loi jointe de $(X_{t_1 + h}, \dots, X_{t_k + h})$ ne
dépend pas de $h$. Il possède une \emph{stationnarité faible}\index{stationnarité faible} si sa moyenne est constante et si $\Cov(X_t, X_{t+h})$ ne dépend que de $h$.
Un \emph{bruit blanc}\index{bruit blanc} est une suite faiblement stationnaire, de moyenne nulle et sans autocorrélation à aucun retard non nul.
\end{definition}

\begin{definition}[Fonctions d'autocovariance, d'autocorrélation et d'autocorrélation partielle]\label{def:qm:linear-time-series:acf}
Pour un processus faiblement stationnaire, la \emph{fonction d'autocovariance}\index{fonction d'autocovariance} est $\gamma(h) = \Cov(X_t, X_{t+h})$, et la
\emph{fonction d'autocorrélation}\index{fonction d'autocorrélation} est $\rho(h) = \gamma(h)/\gamma(0)$. La \emph{fonction d'autocorrélation
partielle}\index{fonction d'autocorrélation partielle} au retard $h$ est le dernier coefficient $\phi_{hh}$ du meilleur prédicteur linéaire de $X_{t+h}$ à partir de $X_{t+h-1},
\dots, X_t$~: la corrélation au retard $h$ une fois pris en compte les retards intermédiaires.
\end{definition}

Sous \omterm{def:qm:linear-time-series:stationary}{bruit blanc}, les autocorrélations empiriques sont approximativement indépendantes, de loi $\mathcal N(0, 1/n)$, d'où les bandes $\pm 2/\sqrt n$ de chaque
corrélogramme. La décomposition qui justifie les modèles de la section suivante est celle de Wold (1938)~: tout processus faiblement stationnaire, purement
non déterministe est une moyenne mobile infinie $X_t = \mu + \sum_{j \ge 0}\psi_j\varepsilon_{t-j}$ de ses erreurs de prédiction à un pas
$\varepsilon_t$, un \omterm{def:qm:linear-time-series:stationary}{bruit blanc}, avec $\psi_0 = 1$ et $\sum_j\psi_j^2 < \infty$.

Le corrélogramme du spread 2s10s (\cref{fig:qm:linear-time-series:acf}) est la signature d'un niveau presque non stationnaire~: des autocorrélations
qui partent de 0.9987 et décroissent presque linéairement sur un an de retards. Ses variations quotidiennes sont presque un \omterm{def:qm:linear-time-series:stationary}{bruit blanc}, avec un écart type de 4.6
points de base, et des autocorrélations de 0.015, 0.015 et $-0.027$ aux trois premiers retards, dans la bande $\pm 0.018$ pour les deux premiers.

\begin{omfigure}
\begin{tikzpicture}
\begin{groupplot}[group style={group size=2 by 1,horizontal sep=1.5cm},width=7.4cm,height=4.8cm,
  tick label style={font=\small},label style={font=\small},grid=major,grid style={gray!25},table/col sep=comma]
\nextgroupplot[xlabel={retard (jours)},ylabel={autocorrélation},xmin=0,xmax=250,ymin=0.6,ymax=1,title={\small niveau du spread}]
\addplot[omDef,very thick] table[x=lag,y=levels]{figdata/methods/17-linear-time-series/acf.csv};
\nextgroupplot[xlabel={retard (jours)},xmin=0,xmax=40,ymin=-0.06,ymax=0.06,title={\small variation quotidienne},
  scaled y ticks=false,yticklabel style={/pgf/number format/fixed,/pgf/number format/precision=2}]
\addplot[ycomb,omDef,thick,mark=*,mark size=1pt] table[x=lag,y=changes]{figdata/methods/17-linear-time-series/acf.csv};
\addplot[black,dashed,sharp plot] coordinates {(0,0.0178) (40,0.0178)};
\addplot[black,dashed,sharp plot] coordinates {(0,-0.0178) (40,-0.0178)};
\end{groupplot}
\end{tikzpicture}

\omcaption{Autocorrélations empiriques du spread 2s10s du Trésor américain, du 1\up{er} juin 1976 au 22 septembre 2026 (12\,574 jours). À gauche~: le niveau, retards 1 à
250, en partant de 0.9987. À droite~: la variation quotidienne, retards 1 à 40, avec la bande $\pm 2/\sqrt n$. Données~: séries FRED DGS10 et DGS2 (Conseil
des gouverneurs de la Réserve fédérale, H.15).}\label{fig:qm:linear-time-series:acf}
\end{omfigure}

\section{Modèles autorégressifs et à moyenne mobile}

\begin{definition}[Processus autorégressifs, à moyenne mobile et ARMA]\label{def:qm:linear-time-series:arma}
Avec $(\varepsilon_t)$ un \omterm{def:qm:linear-time-series:stationary}{bruit blanc} de variance $\sigma^2$~: un \emph{processus autorégressif}\index{processus autorégressif} d'ordre $p$, AR($p$), vérifie
$X_t = c + \sum_{i=1}^p\phi_iX_{t-i} + \varepsilon_t$~; un \emph{processus moyenne mobile}\index{processus moyenne mobile} d'ordre $q$, MA($q$), est $X_t = \mu +
\varepsilon_t + \sum_{j=1}^q\vartheta_j\varepsilon_{t-j}$~; un \emph{processus ARMA}\index{processus ARMA} les combine, $\phi(L)X_t = c + \vartheta(L)\varepsilon_t$ avec
$\phi(z) = 1 - \sum_i\phi_iz^i$, $\vartheta(z) = 1 + \sum_j\vartheta_jz^j$ et $L$ l'opérateur de retard.
\end{definition}

\begin{proposition}[Stationnarité et inversibilité d'un processus ARMA]\label{prop:qm:linear-time-series:roots}
L'équation ARMA admet une unique solution faiblement stationnaire et causale $X_t = \mu + \sum_{j \ge 0}\psi_j\varepsilon_{t-j}$ si et seulement si toute racine de $\phi(z)$
se trouve à l'extérieur du cercle unité~; le processus est inversible, $\varepsilon_t = \sum_{j \ge 0}\pi_jX_{t-j}$ à une constante près, si et seulement si toute racine de $\vartheta(z)$
le fait, les polynômes n'ayant aucune racine commune.
\end{proposition}

\begin{proof}
On factorise $\phi(z) = \prod_i(1 - z/z_i)$. Pour $|z_i| > 1$, $(1 - L/z_i)^{-1} = \sum_kz_i^{-k}L^k$ est un filtre absolument sommable, donc $\psi(z) = \vartheta(z)/\phi(z)$
a des coefficients sommables et définit la solution~; si certains $|z_i| \le 1$, aucun inverse causal sommable n'existe (pour $|z_i| = 1$, il n'en existe pas du tout). Le cas MA
est le même raisonnement appliqué à $\vartheta$.
\end{proof}

Pour l'AR(1) $X_t = c + \phi X_{t-1} + \varepsilon_t$, la condition est $|\phi| < 1$, les autocorrélations valent $\phi^h$, et un écart à la moyenne
se réduit de moitié en $\ln(1/2)/\ln\phi$ périodes, soit la \omterm{def:qm:stochastic-differential-equations:ou}{demi-vie} du \omterm{def:qm:stochastic-differential-equations:ou}{processus d'Ornstein--Uhlenbeck} du chapitre~4, échantillonné quotidiennement. Un AR($p$) possède une \omterm{def:qm:linear-time-series:acf}{fonction d'autocorrélation
partielle} qui s'annule après le retard $p$~; un MA($q$) possède une \omterm{def:qm:linear-time-series:acf}{fonction d'autocorrélation} qui s'annule après le retard $q$. Les ordres sont
choisis par un \omterm{def:qm:linear-time-series:ic}{critère d'information}.

\begin{definition}[Critère d'information]\label{def:qm:linear-time-series:ic}
Un \emph{critère d'information}\index{critère d'information} classe les modèles selon leur ajustement pénalisé par leur taille~: celui d'Akaike, $\mathrm{AIC} = -2\ell + 2k$, et celui de Schwarz
$\mathrm{BIC} = -2\ell + k\ln n$, pour une log-vraisemblance maximisée $\ell$ et $k$ paramètres~; le plus petit l'emporte.
\end{definition}

L'AIC vise la prévision et tend à retenir de grands modèles~; le BIC est convergent pour l'ordre vrai lorsqu'il existe. Sur les variations quotidiennes du 2s10s, le critère AIC
retient un AR(10) dont tous les coefficients sont inférieurs à 0.035 en valeur absolue~: statistiquement détectables avec 12\,574 jours, économiquement nuls.

\section{La vision spectrale}

\begin{definition}[Densité spectrale, périodogramme]\label{def:qm:linear-time-series:spectral}
La \emph{densité spectrale}\index{densité spectrale} d'un processus faiblement stationnaire à autocovariances sommables est $f(\omega) = \frac1{2\pi}\sum_h\gamma(h)e^{-ih\omega}$,
$\omega \in [-\pi, \pi]$~: la décomposition de sa variance par fréquence, $\gamma(0) = \int_{-\pi}^\pi f$. Le \emph{périodogramme}\index{périodogramme} $I(\omega_j) =
\frac1{2\pi n}\bigl\lvert\sum_tX_te^{-it\omega_j}\bigr\rvert^2$ aux fréquences de Fourier $\omega_j = 2\pi j/n$ en est l'analogue empirique.
\end{definition}

Le \omterm{def:qm:linear-time-series:stationary}{bruit blanc} a un spectre plat, $\sigma^2/2\pi$~; un AR(1) avec $\phi$ proche de un concentre sa variance aux basses fréquences, $f(\omega) =
\sigma^2/(2\pi|1 - \phi e^{-i\omega}|^2)$. Le \omterm{def:qm:linear-time-series:spectral}{périodogramme} est asymptotiquement sans biais mais non convergent~: chaque ordonnée vaut approximativement $f(\omega_j)$ multipliée par
une variable exponentielle, quelle que soit la taille de l'échantillon, si bien qu'il faut le lisser ou le moyenner. Sa partie basse fréquence est celle où la persistance se manifeste, et elle est
le fondement de l'estimation de la \omterm{def:qm:linear-time-series:longmem}{mémoire longue} de la dernière section. $f(0) = \frac1{2\pi}\sum_h\gamma(h)$ est, à $2\pi$ près, la variance de long terme du
chapitre~11.

\section{Racines unitaires}

\begin{definition}[Racine unitaire, régression fallacieuse]\label{def:qm:linear-time-series:unitroot}
Un processus possède une \emph{racine unitaire}\index{racine unitaire} si son polynôme autorégressif admet une racine en $z = 1$~: sa différence première est stationnaire, mais le niveau
n'est pas stationnaire, comme une marche aléatoire. Une \emph{régression fallacieuse}\index{régression fallacieuse} est une régression entre processus indépendants à racines unitaires dont les statistiques
$t$ usuelles et le $R^2$ signalent une relation qui n'existe pas.
\end{definition}

Granger et Newbold (1974) ont mis le problème en évidence par simulation, et il est facile de le reproduire~: en régressant une marche aléatoire de 500 pas sur une autre,
indépendante, la statistique $t$ de la pente dépasse 1.96 en valeur absolue dans 89~\% des 5\,000 essais, avec un $R^2$ médian de 0.17. Le résidu
hérite de la \omterm{def:qm:linear-time-series:unitroot}{racine unitaire}, les observations ne sont pas indépendantes, et toute formule qui les supposait indépendantes est fausse. La même défaillance touche le
test de $\phi = 1$ lui-même.

\begin{definition}[Test de Dickey--Fuller]\label{def:qm:linear-time-series:df}
Le \emph{test de Dickey--Fuller}\index{test de Dickey--Fuller} d'une \omterm{def:qm:linear-time-series:unitroot}{racine unitaire} régresse $\Delta X_t$ sur $X_{t-1}$ (avec une constante, et une tendance si on le souhaite) et
compare la statistique $t$ $\tau$ du coefficient de $X_{t-1}$ à la loi de Dickey--Fuller plutôt qu'à la loi normale~; la version augmentée
ajoute des différences retardées $\Delta X_{t-k}$ pour absorber la dépendance de court terme (Said et Dickey, 1984).
\end{definition}

Sous l'hypothèse de \omterm{def:qm:linear-time-series:unitroot}{racine unitaire}, avec une constante, $\tau$ converge vers $\bigl(\frac12(W(1)^2 - 1) - W(1)\int_0^1W\bigr)/\bigl(\int_0^1W^2 - (\int_0^1W)^2\bigr)^{1/2}$ pour un
\omterm{def:qm:brownian-motion:bm}{mouvement brownien} standard $W$, une fonctionnelle qui n'est ni normale ni centrée en zéro (\cref{fig:qm:linear-time-series:df}). Ses quantiles n'ont pas de
forme close~; les surfaces de réponse de MacKinnon (2010) donnent les valeurs critiques à 1~\%, 5~\% et 10~\% $-3.43$, $-2.86$ et $-2.57$ pour de grands échantillons. La simulation propre au chapitre
de 20\,000 marches aléatoires de 500 pas donne $-3.46$, $-2.89$ et $-2.57$, contre $-3.44$, $-2.87$ et $-2.57$ chez MacKinnon à cette longueur.

\begin{omfigure}
\begin{tikzpicture}
\begin{axis}[width=10cm,height=5cm,xlabel={statistique $t$ du coefficient de $X_{t-1}$},ylabel={densité},
  tick label style={font=\small},label style={font=\small},xmin=-5,xmax=3,ymin=0,ymax=0.55,grid=major,grid style={gray!25},
  legend columns=2,legend style={font=\footnotesize,at={(0.5,-0.3)},anchor=north,draw=none,fill=none,
  /tikz/every even column/.append style={column sep=8pt}},table/col sep=comma]
\addplot[omDef,thick,const plot mark mid,fill=omDef!15] table[x=tau,y=density]{figdata/methods/17-linear-time-series/df.csv};
\addplot[omThm,very thick,dashed] table[x=tau,y=normal]{figdata/methods/17-linear-time-series/df.csv};
\addplot[black,thin,sharp plot] coordinates {(-2.86,0) (-2.86,0.55)};
\addplot[black,thin,dotted,sharp plot] coordinates {(-1.645,0) (-1.645,0.55)};
\legend{Dickey--Fuller (simulée),loi normale}
\end{axis}
\end{tikzpicture}

\omcaption{La loi de la statistique $t$ de Dickey--Fuller (régression avec constante) sous \omterm{def:qm:linear-time-series:unitroot}{racine unitaire}, à partir de 20\,000 marches aléatoires simulées de 500 pas,
face à la loi normale centrée réduite. Les valeurs critiques à 5~\% sont $-2.86$ (trait plein) et $-1.645$ (pointillés)~: un test fondé sur la table normale rejette à tort une \omterm{def:qm:linear-time-series:unitroot}{racine unitaire}
vraie beaucoup trop souvent. Données~: le tutoriel du chapitre, avec graine fixée.}\label{fig:qm:linear-time-series:df}
\end{omfigure}

Pour le spread 2s10s, $\tau = -2.79$ sans différence retardée~: non significatif à 5~\%, significatif à 10~\%. Le test augmenté dépend des retards~:
$-2.74$ avec cinq retards, $-3.13$ avec dix, $-3.55$ avec vingt~; l'AIC, en cherchant jusqu'à quarante, retient 34 et $-3.14$. Sur l'échantillon depuis 1990, la statistique de Dickey--Fuller
vaut $-2.01$. La lecture honnête est que les données ne peuvent pas trancher~: le spread revient vers sa moyenne trop lentement pour que cinquante ans de données quotidiennes permettent de le distinguer
d'une marche aléatoire.

\begin{proposition}[Biais de Kendall du coefficient AR(1)]\label{prop:qm:linear-time-series:kendall}
Pour un AR(1) stationnaire à moyenne estimée, le coefficient des moindres carrés vérifie $\E[\hat\phi] - \phi = -(1 + 3\phi)/n + O(n^{-2})$ (Kendall, 1954).
\end{proposition}

Le biais est faible sur le coefficient et grand sur la \omterm{def:qm:stochastic-differential-equations:ou}{demi-vie}, car la \omterm{def:qm:stochastic-differential-equations:ou}{demi-vie} divise par $\ln\phi$, proche de zéro. L'estimation du spread
0.99875, corrigée de $(1 + 3 \times 0.99875)/12\,574 = 0.00032$, devient 0.99906, et la \omterm{def:qm:stochastic-differential-equations:ou}{demi-vie} passe de 552 jours de cotation (2.2 ans) à 739 (2.9
ans). La correction n'est qu'une moyenne~: en simulant 4\,000 historiques de 12\,574 jours à partir d'un AR(1) de $\phi = 0.99906$, la \omterm{def:qm:stochastic-differential-equations:ou}{demi-vie} estimée a pour médiane
576 jours et un intervalle 10--90~\% de 352 à 969 jours (\cref{fig:qm:linear-time-series:halflife}), et dans 57~\% des historiques, le \omterm{def:qm:linear-time-series:df}{test de Dickey--Fuller}
ne rejette pas la \omterm{def:qm:linear-time-series:unitroot}{racine unitaire} au seuil de 5~\%, alors qu'il n'y en a aucune.

\begin{omfigure}
\begin{tikzpicture}
\begin{axis}[width=10cm,height=5cm,xlabel={demi-vie estimée (jours de cotation)},ylabel={densité ($\times 10^{-3}$)},
  tick label style={font=\small},label style={font=\small},xmin=0,xmax=2000,ymin=0,ymax=2.2,grid=major,grid style={gray!25},
  xtick={0,500,1000,1500,2000},xticklabel style={/pgf/number format/set thousands separator={}},
  legend style={font=\footnotesize,at={(0.98,0.96)},anchor=north east,fill=white,draw=none},table/col sep=comma]
\addplot[omDef,thick,const plot mark mid,fill=omDef!15] table[x=days,y=density]{figdata/methods/17-linear-time-series/halflife.csv};
\addplot[black,dashed,thick,sharp plot] coordinates {(739,0) (739,2.2)};
\addplot[omThm,very thick,sharp plot] coordinates {(552,0) (552,2.2)};
\legend{estimations simulées,demi-vie vraie (739),estimation 2s10s (552)}
\end{axis}
\end{tikzpicture}

\omcaption{Distribution d'échantillonnage de la \omterm{def:qm:stochastic-differential-equations:ou}{demi-vie} de l'AR(1) estimée à partir de 12\,574 observations quotidiennes d'un AR(1) dont la \omterm{def:qm:stochastic-differential-equations:ou}{demi-vie} vraie est de 739 jours, sur
4\,000 historiques simulés (les estimations supérieures à 2\,000 jours sont regroupées dans la dernière classe)~: médiane 576, intervalle 10--90~\% de 352 à 969. L'estimation 2s10s de 552 jours est un
tirage typique. Données~: le tutoriel du chapitre, avec graine fixée.}\label{fig:qm:linear-time-series:halflife}
\end{omfigure}

\section{Mémoire longue}

\begin{definition}[Mémoire longue, différenciation fractionnaire, exposant de Hurst]\label{def:qm:linear-time-series:longmem}
Un \omterm{def:qm:linear-time-series:stationary}{processus stationnaire} possède une \emph{mémoire longue}\index{mémoire longue} si ses autocorrélations décroissent comme une puissance, $\rho(h) \sim ch^{2d-1}$ avec $0 < d < \frac12$, de sorte
qu'elles ne soient pas sommables. La \emph{différenciation fractionnaire}\index{différenciation fractionnaire} applique $(1 - L)^d = \sum_{j \ge 0}w_jL^j$ avec $w_0 = 1$, $w_j = w_{j-1}(j - 1 -
d)/j$~; le processus ARFIMA$(0, d, 0)$ est $(1 - L)^dX_t = \varepsilon_t$ (Granger et Joyeux, 1980~; Hosking, 1981). L'\emph{exposant de Hurst}\index{exposant de Hurst} est
$H = d + \frac12$, mesuré par la croissance des étendues ou des variances des sommes partielles, comme $n^H$ (Hurst, 1951).
\end{definition}

\begin{definition}[Mouvement brownien fractionnaire]\label{def:qm:linear-time-series:fbm}
Un \emph{mouvement brownien fractionnaire}\index{mouvement brownien fractionnaire} d'\omterm{def:qm:linear-time-series:longmem}{exposant de Hurst} $H \in (0, 1)$ est le \omterm{def:qm:brownian-motion:bm}{processus gaussien} centré tel que $B_H(0) = 0$ et
$\E[B_H(t)B_H(s)] = \frac12(t^{2H} + s^{2H} - |t - s|^{2H})$ (Mandelbrot et Van Ness, 1968). Pour $H = \frac12$, c'est le \omterm{def:qm:brownian-motion:bm}{mouvement brownien}~; pour $H > \frac12$, ses accroissements
sont corrélés positivement, avec une décroissance en loi de puissance, et pour $H < \frac12$ négativement.
\end{definition}

Pour $H \ne \frac12$, le \omterm{def:qm:linear-time-series:fbm}{mouvement brownien fractionnaire} n'est pas une \omterm{def:qm:jump-processes:jd}{semimartingale}, de sorte que le calcul d'Itô du chapitre~3 ne s'y applique pas, et un prix qui le suivrait
admettrait un arbitrage~; c'est un modèle de volatilité et de quantités à retour lent à la moyenne, non pour les prix. La différence entre \omterm{def:qm:linear-time-series:longmem}{mémoire longue} et mémoire
courte tient à la queue du corrélogramme (\cref{fig:qm:linear-time-series:longmem})~: un ARFIMA avec $d = 0.3$ et un AR(1) ayant la même première
autocorrélation, 0.41, diffèrent d'un facteur mille au retard 10. Sur 20\,000 observations simulées, l'estimation de $H$ par l'étendue réévaluée donne 0.78 et
l'estimation par \omterm{def:qm:linear-time-series:spectral}{log-périodogramme} (Geweke--Porter-Hudak) donne 0.71, pour une vraie valeur de 0.8~; sur un \omterm{def:qm:linear-time-series:stationary}{bruit blanc}, elles donnent 0.53 et 0.46, l'étendue réévaluée étant biaisée à la hausse dans
les blocs courts.

\begin{omfigure}
\begin{tikzpicture}
\begin{loglogaxis}[width=10cm,height=5cm,xlabel={retard},ylabel={autocorrélation},
  tick label style={font=\small},label style={font=\small},xmin=1,xmax=200,ymin=1e-4,ymax=1,grid=major,grid style={gray!25},
  legend columns=3,legend style={font=\footnotesize,at={(0.5,-0.3)},anchor=north,draw=none,fill=none,
  /tikz/every even column/.append style={column sep=8pt}},table/col sep=comma]
\addplot[only marks,mark=*,mark size=0.9pt,omDef] table[x=lag,y=arfima]{figdata/methods/17-linear-time-series/longmem.csv};
\addplot[omDef,very thick] table[x=lag,y=theory]{figdata/methods/17-linear-time-series/longmem.csv};
\addplot[omThm,very thick,dashed] table[x=lag,y=ar1]{figdata/methods/17-linear-time-series/longmem.csv};
\legend{{ARFIMA, $d = 0.3$, échantillon},{ARFIMA, théorie},{AR(1), même $\rho(1)$}}
\end{loglogaxis}
\end{tikzpicture}

\omcaption{Autocorrélations d'un processus à \omterm{def:qm:linear-time-series:longmem}{mémoire longue} (ARFIMA avec $d = 0.3$~: 20\,000 observations simulées et la valeur exacte $\rho(h) = \prod_{i \le h}(i - 1 + d)/(i -
d)$, une droite de pente $2d - 1$ sur ces axes) et d'un AR(1) de même autocorrélation de retard un, qui décroît géométriquement. Données~: le
tutoriel du chapitre, avec graine fixée.}\label{fig:qm:linear-time-series:longmem}
\end{omfigure}

Les variations quotidiennes du spread donnent une estimation par \omterm{def:qm:linear-time-series:spectral}{log-périodogramme} de $H = 0.34$ à partir des $\sqrt n = 112$ fréquences les plus basses, soit $d = -0.16$ avec une \omterm{def:qm:estimation:estimator}{erreur type}
de $\pi/\sqrt{24 \times 112} = 0.06$~: les variations sont surdifférenciées, et le niveau se comporte comme un processus intégré fractionnairement de $d \approx 0.84$, à \omterm{def:qm:stochastic-differential-equations:ou}{retour
à la moyenne} à long terme, mais avec une décroissance hyperbolique plutôt que géométrique. C'est l'une des raisons pour lesquelles les \omterm{def:qm:stochastic-differential-equations:ou}{demi-vies} AR(1) du spread sont instables
d'un échantillon à l'autre~: l'AR(1) n'a pas la bonne forme de mémoire.

\section{Tutoriel~: la demi-vie d'un spread}\label{tut:qm:linear-time-series:halflife}

\begin{tutorial}
\textbf{Objectif.} Estimer la persistance du spread 2s10s de trois façons et décider ce que les données permettent de dire sur une \omterm{def:qm:linear-time-series:unitroot}{racine unitaire}.
\textbf{État final~:} \cref{fig:qm:linear-time-series:acf,fig:qm:linear-time-series:df,fig:qm:linear-time-series:halflife}~; les \omterm{def:qm:stochastic-differential-equations:ou}{demi-vies} de 552 et 739 jours et les
statistiques de Dickey--Fuller.
\begin{tutsteps}
\item \textbf{L'AR(1) avec sa \omterm{def:qm:stochastic-differential-equations:ou}{demi-vie} et la correction de Kendall.}
\omcode{code/firm/tsa/firm_tsa.py}{102}{118}{AR(1) par moindres carrés, demi-vie et correction de biais.}\label{lst:qm:linear-time-series:ar1}
\item \textbf{La régression de Dickey--Fuller augmentée et les valeurs critiques de MacKinnon.}
\omcode{code/firm/tsa/firm_tsa.py}{125}{158}{Le test de Dickey--Fuller augmenté.}\label{lst:qm:linear-time-series:adf}
\item \textbf{Exécuter} \texttt{problem()}, \texttt{half\_life\_mc(0.99906, 12574)}, \texttt{spurious()} et \texttt{long\_memory()} dans \texttt{qm\_tsa.py}, puis
\texttt{fig\_tsa.py}.
\end{tutsteps}
\textbf{Pour aller plus loin.} Ajuster un ARFIMA au spread par le $d$ du \omterm{def:qm:linear-time-series:spectral}{log-périodogramme} et comparer sa prévision de la décroissance du spread à celle de l'AR(1)~; relancer les
tests sur des données hebdomadaires et voir quelles conclusions tiennent.
\end{tutorial}

\section{Construction~: la boîte à outils des séries temporelles}\label{bld:qm:linear-time-series:tsa}

\begin{build}
\bfield{Rôle} Toute persistance sur laquelle s'appuie la mini-firme (le \omterm{def:qm:stochastic-differential-equations:ou}{retour à la moyenne} d'un spread, la décroissance d'un signal, la mémoire d'une volatilité) est mesurée et testée
ici avant qu'une stratégie ne soit bâtie dessus.
\bfield{Interface} \texttt{acf}, \texttt{pacf}, \texttt{band}~; \texttt{yule\_walker(x, p)}~; \texttt{arma\_css(x, p, q)}~; \texttt{ar1(x)}~; \texttt{adf(x, lags,
trend, max\_lags)}~; \texttt{mackinnon\_crit(trend, n)}~; \texttt{periodogram(x)}~; \texttt{frac\_diff(x, d, k)}~; \texttt{hurst\_rs(x)}, \texttt{hurst\_gph(x)}.
\bfield{Règles} Les tests de \omterm{def:qm:linear-time-series:unitroot}{racine unitaire} signalent leur choix de retards et sont relancés sur une plage de retards~; les \omterm{def:qm:stochastic-differential-equations:ou}{demi-vies} sont données avec la correction de biais et un
intervalle simulé~; aucune table normale n'est utilisée pour un coefficient proche de un.
\bfield{Tests d'acceptation} \texttt{code/firm/tsa/tests/}~: autocorrélations et autocorrélations partielles d'un AR(1)~; Yule--Walker pour un AR(2)~; ARMA(1, 1) par
moindres carrés conditionnels~; le \omterm{def:qm:linear-time-series:df}{test de Dickey--Fuller} a sa taille nominale sur les marches aléatoires et rejette un AR(1) stationnaire~; la formule de MacKinnon~; un \omterm{def:qm:linear-time-series:spectral}{périodogramme} plat pour
un \omterm{def:qm:linear-time-series:stationary}{bruit blanc}~; la \omterm{def:qm:linear-time-series:longmem}{différenciation fractionnaire} avec $d = 1$ est une différenciation~; les estimations de Hurst d'un \omterm{def:qm:linear-time-series:stationary}{bruit blanc} et d'un ARFIMA.
\bfield{Pour aller plus loin} Le test KPSS, dont l'\omterm{def:qm:testing-and-multiple-testing:test}{hypothèse nulle} est la stationnarité~; la vraisemblance gaussienne exacte d'un ARMA par le filtre de Kalman du chapitre~19~; l'\omterm{def:qm:estimation:estimator}{estimateur} de Whittle local de $d$.
\end{build}

\begin{omsources}
\item H.~Wold, \emph{A Study in the Analysis of Stationary Time Series}, Almqvist and Wiksell, 1938.
\item D.~A.~Dickey and W.~A.~Fuller, ``Distribution of the estimators for autoregressive time series with a unit root'', \emph{Journal of the American Statistical
Association} 74, 1979~; S.~E.~Said and D.~A.~Dickey, \emph{Biometrika} 71, 1984.
\item J.~G.~MacKinnon, ``Critical values for cointegration tests'', Queen's Economics Department Working Paper 1227, 2010.
\item C.~W.~J.~Granger and P.~Newbold, ``Spurious regressions in econometrics'', \emph{Journal of Econometrics} 2, 1974.
\item M.~G.~Kendall, ``Note on bias in the estimation of autocorrelation'', \emph{Biometrika} 41, 1954.
\item H.~E.~Hurst, ``Long-term storage capacity of reservoirs'', \emph{Transactions of the American Society of Civil Engineers} 116, 1951.
\item B.~B.~Mandelbrot and J.~W.~Van Ness, ``Fractional Brownian motions, fractional noises and applications'', \emph{SIAM Review} 10, 1968.
\item C.~W.~J.~Granger and R.~Joyeux, \emph{Journal of Time Series Analysis} 1, 1980~; J.~R.~M.~Hosking, ``Fractional differencing'', \emph{Biometrika} 68, 1981~;
J.~Geweke and S.~Porter-Hudak, \emph{Journal of Time Series Analysis} 4, 1983.
\item Board of Governors of the Federal Reserve System, H.15 Selected Interest Rates, via FRED (series DGS2 and DGS10), consulté le 24 septembre 2026.
\end{omsources}

\section{Exercices}

\begin{exercise}[$\star$]\label{exo:qm:linear-time-series:1}
Un AR(1) a un coefficient de 0.98 sur des données quotidiennes. Quelle est la \omterm{def:qm:stochastic-differential-equations:ou}{demi-vie} en jours, et quelle est l'autocorrélation de retard 20~?
\end{exercise}

\begin{exercise}[$\star$]\label{exo:qm:linear-time-series:2}
Le MA(1) $X_t = \varepsilon_t + 2\varepsilon_{t-1}$ est-il inversible~? Trouver le MA(1) inversible ayant les mêmes autocorrélations.
\end{exercise}

\begin{exercise}[$\star$]\label{exo:qm:linear-time-series:3}
Écrire les quatre premiers poids de \omterm{def:qm:linear-time-series:longmem}{différenciation fractionnaire} pour $d = 0.3$.
\end{exercise}

\begin{exercise}[$\star\star$]\label{exo:qm:linear-time-series:4}
Pour un AR(1) avec $\phi = 0.99$ estimé sur 1\,000 observations quotidiennes, quel est le biais de Kendall, et de combien modifie-t-il la \omterm{def:qm:stochastic-differential-equations:ou}{demi-vie}~?
\end{exercise}

\begin{exercise}[$\star\star$]\label{exo:qm:linear-time-series:5}
Montrer que les autocorrélations d'un AR(2) vérifient $\rho(h) = \phi_1\rho(h - 1) + \phi_2\rho(h - 2)$, puis calculer $\rho(1)$ et $\rho(2)$ pour $\phi_1 = 0.5$, $\phi_2 =
-0.3$.
\end{exercise}

\begin{exercise}[$\star\star$]\label{exo:qm:linear-time-series:6}
En utilisant les coefficients de MacKinnon $(-2.86154, -2.8903, -4.234, -40.040)$, calculer la valeur critique de Dickey--Fuller à 5~\% (avec constante) pour $T = 100$ et $T = 12\,574$.
\end{exercise}

\begin{exercise}[$\star\star\star$]\label{exo:qm:linear-time-series:7}
\emph{Programmation.} Régresser une marche aléatoire indépendante de 500 pas sur une autre, 5\,000 fois. À quelle fréquence a-t-on $|t| > 1.96$~? Que se passe-t-il si l'on régresse les différences
plutôt~?
\end{exercise}

\begin{exercise}[$\star\star\star$]\label{exo:qm:linear-time-series:8}
\emph{Trouver la faille.} « Nous avons régressé la variation quotidienne du spread sur le niveau retardé~; le coefficient vaut $-0.00125$ avec $t = -2.79$, significatif à 1~\%, donc le
spread revient à la moyenne et nous allons trader son retour avec une \omterm{def:qm:stochastic-differential-equations:ou}{demi-vie} de 552 jours. »
\end{exercise}

\section{Problème~: la demi-vie d'un spread}

\begin{problem}[{Problème du week-end --- à quelle vitesse le spread 2s10s revient-il à la moyenne~?}]\label{pb:qm:linear-time-series:1}
Les données sont les rendements quotidiens du Trésor à dix ans et à deux ans à échéance constante (FRED DGS10 et DGS2) du 1\up{er} juin 1976 au 22 septembre 2026, sur les 12\,574
jours où les deux sont disponibles~; le spread est leur différence en points de base.

\medskip\noindent\textbf{Partie I --- La série.}
\begin{enumerate}
\item Quelles sont la moyenne, l'étendue et la dernière valeur du spread~?
\item Quelle est son autocorrélation d'ordre un, et à quoi ressemble son corrélogramme~?
\item Quels sont l'écart type et les premières autocorrélations de ses variations quotidiennes~?
\item Quel ordre AR le critère AIC retient-il pour les variations, et cela importe-t-il économiquement~?
\item Quelle est, approximativement, la représentation de Wold des variations~?
\end{enumerate}

\medskip\noindent\textbf{Partie II --- L'AR(1).}
\begin{enumerate}[resume]
\item Quels sont le coefficient de l'AR(1), son \omterm{def:qm:estimation:estimator}{erreur type} et la \omterm{def:qm:stochastic-differential-equations:ou}{demi-vie}~?
\item Quelle est la statistique $t$ par rapport à un, et quelle $p$-valeur donne la table normale~?
\item Quelle est la correction de Kendall, et quelle est la \omterm{def:qm:stochastic-differential-equations:ou}{demi-vie} corrigée~?
\item Que disent les \omterm{def:qm:linear-time-series:df}{tests de Dickey--Fuller} et de Dickey--Fuller augmenté, avec 0, 5, 10 et 20 retards~?
\item Que donne l'échantillon depuis 1990~?
\end{enumerate}

\medskip\noindent\textbf{Partie III --- Ce que les données peuvent dire.}
\begin{enumerate}[resume]
\item Si la véritable \omterm{def:qm:stochastic-differential-equations:ou}{demi-vie} était de 739 jours, comment se répartiraient les estimations sur 12\,574 jours~?
\item À quelle fréquence le \omterm{def:qm:linear-time-series:df}{test de Dickey--Fuller} ne rejetterait-il alors pas la \omterm{def:qm:linear-time-series:unitroot}{racine unitaire}~?
\item Que dit le \omterm{def:qm:linear-time-series:spectral}{log-périodogramme} des variations, et du niveau~?
\item Pourquoi une \omterm{def:qm:stochastic-differential-equations:ou}{demi-vie} d'AR(1) pourrait-elle être instable d'un échantillon à l'autre~?
\item Que voudrait savoir un trader en steepener que la \omterm{def:qm:stochastic-differential-equations:ou}{demi-vie} ne dit pas~?
\end{enumerate}

\medskip\noindent\textbf{Partie IV --- Jugement.}
\begin{enumerate}[resume]
\item Le spread est-il stationnaire~?
\item Quelle \omterm{def:qm:stochastic-differential-equations:ou}{demi-vie} utiliseriez-vous pour dimensionner une transaction de \omterm{def:qm:stochastic-differential-equations:ou}{retour à la moyenne}, et avec quelle réserve~?
\item Pourquoi la $p$-valeur de 0.003 fournie par la table normale est-elle trompeuse~?
\item Énoncer le \emph{résultat nommé}~: la \omterm{def:qm:stochastic-differential-equations:ou}{demi-vie} de l'AR(1), sa valeur corrigée du biais, et si une \omterm{def:qm:linear-time-series:unitroot}{racine unitaire} peut être rejetée.
\item En une phrase~: que nous apprend un test de \omterm{def:qm:linear-time-series:unitroot}{racine unitaire} pour un trader~?
\end{enumerate}
\end{problem}

\section{Questions d'entretien}

\begin{interviewq}[$\star$ \iqroles{chercheur, trader}]\label{iq:qm:linear-time-series:1}
Quelle est la \omterm{def:qm:stochastic-differential-equations:ou}{demi-vie} d'un AR(1) de coefficient $\phi$, et pourquoi est-elle difficile à estimer lorsque $\phi$ est proche de un~?
\end{interviewq}

\begin{interviewq}[$\star\star$ \iqroles{chercheur}]\label{iq:qm:linear-time-series:2}
Pourquoi ne peut-on pas utiliser la table $t$ usuelle pour tester $\phi = 1$~?
\end{interviewq}

\begin{interviewq}[$\star\star$ \iqroles{chercheur, ingénieur ML}]\label{iq:qm:linear-time-series:3}
Vous régressez une série de prix sur une autre et obtenez $R^2 = 0.6$ et $t = 15$. Que vérifiez-vous~?
\end{interviewq}

\begin{interviewq}[$\star\star$ \iqroles{chercheur}]\label{iq:qm:linear-time-series:4}
Comment distinguez-vous un processus AR d'un processus MA à partir de leurs corrélogrammes~?
\end{interviewq}

\begin{interviewq}[$\star\star$ \iqroles{chercheur, ingénieur ML}]\label{iq:qm:linear-time-series:5}
Qu'est-ce que la \omterm{def:qm:linear-time-series:longmem}{mémoire longue}, et où l'observez-vous sur les marchés~?
\end{interviewq}

\begin{interviewq}[$\star\star\star$ \iqroles{chercheur}]\label{iq:qm:linear-time-series:6}
Pourquoi le coefficient AR(1) estimé par moindres carrés est-il biaisé à la baisse sur petit échantillon, et de quelle ampleur est ce biais~?
\end{interviewq}
